\documentclass[a4paper,11pt]{article}

% ---------- Кодировка и язык ----------
\usepackage[utf8]{inputenc}
\usepackage[T2A]{fontenc}
\usepackage[english,russian]{babel}

% ---------- Геометрия страницы ----------
\usepackage[margin=2.5cm]{geometry}

% ---------- Математика ----------
\usepackage{amsmath,amssymb,amsthm,mathtools}
\usepackage{empheq}
\usepackage{bm}          % жирные символы в формулах: \bm{x}

% ---------- Картинки ----------
\usepackage{graphicx}
\graphicspath{{figures/}}   % картинки лежат в latex/figures/
\usepackage{float}          % позиция [H] для картинок
\usepackage{subcaption}     % несколько картинок рядом

% ---------- Оформление ----------
\usepackage{enumitem}
\usepackage{booktabs}       % красивые таблицы
\usepackage{xcolor}
\usepackage{xr}             % ссылки на главы 1--5 из notes.tex (номера берутся из build/notes.aux)
\externaldocument{build/notes}
\usepackage[colorlinks=true,linkcolor=blue,urlcolor=blue,citecolor=blue]{hyperref}
\usepackage{cleveref}       % \cref{...} для умных ссылок
\newcommand{\lean}[1]{\mbox{\path{#1}}} % имена из Lean-кода, без переносов внутри имени (пакет url)

% ---------- Теоремы и окружения ----------
\theoremstyle{plain}
\newtheorem{theorem}{Теорема}[section]
\newtheorem{lemma}[theorem]{Лемма}
\newtheorem{proposition}[theorem]{Утверждение}
\newtheorem{corollary}[theorem]{Следствие}

\theoremstyle{definition}
\newtheorem{definition}[theorem]{Определение}
\newtheorem{example}[theorem]{Пример}

\theoremstyle{remark}
\newtheorem*{remark}{Замечание}

% ---------- Полезные макросы ----------
\newcommand{\R}{\mathbb{R}}
\newcommand{\Z}{\mathbb{Z}}
\newcommand{\N}{\mathbb{N}}
\newcommand{\Cx}{\mathbb{C}}   % \C занят babel-russian
\newcommand{\norm}[1]{\left\lVert #1 \right\rVert}
\newcommand{\abs}[1]{\left\lvert #1 \right\rvert}
\newcommand{\set}[1]{\left\{ #1 \right\}}
\newcommand{\ip}[2]{\langle #1, #2\rangle} % скалярное произведение
\newcommand{\Rpp}[1][2]{\R^{#1}_{++}} % \Rpp = R^2_{++}, \Rpp[n] = R^n_{++}
\newcommand{\Rp}[1][2]{\R^{#1}_{+}}   % \Rp = R^2_{+}, \Rp[n] = R^n_{+}

\DeclareMathOperator{\conv}{conv}
\DeclareMathOperator{\cone}{cone}
\DeclareMathOperator{\Sp}{Sp}

% Заметка на полях / TODO
\newcommand{\todo}[1]{\textcolor{red}{\textbf{TODO:} #1}}

\begin{document}
\setcounter{section}{5} % главы 1--5 и все ссылки на них~--- в notes.tex

\section{Минимальное возмущение цен}

\begin{definition}[Возмущение цен]\label{def:perturbation}
    Для $p,q>0$ назовём \textit{возмущением} величину
    \[
        \pi(p,q):=\max\set{\abs{q/p-1},\ \abs{p/q-1}}.
    \]
    Для наборов цен $P:=\set{p_t\in\Rpp[d]}_{t\in[T]}$ и $Q:=\set{q_t\in\Rpp[d]}_{t\in[T]}$ положим $\pi(P,Q):=\sum_{t\in[T],\,i\in[d]}\pi(p^i_t,q^i_t)$. Для вектора выпусков $y\in\Rp[T]$ назовём \textit{минимальным возмущением} набора цен $P$ величину
    \[
        \rho(P,y):=\inf\set{\pi(P,Q)\mid Q\in\Rpp[d\times T],\ y\ \text{разрешим при ценах}\ Q}.
    \]
    Это наименьшее относительное изменение цен, при котором $y$ становится разрешим. Возмущение симметрично, $\pi(p,q)=\pi(q,p)$, и $\pi(p,q)\to\infty$ при $q\to0$.
\end{definition}

\begin{remark}
    По лемме~\ref{lem:open} множество наборов цен, при которых $y$ разрешим, открыто. Поэтому при $\rho(P,y)>0$ точная нижняя грань не достигается: у допустимого набора $Q$ есть целая окрестность допустимых, и в ней найдётся $Q'$ с меньшим возмущением. По теореме~\ref{thm:closure} замыкание этого множества есть множество наборов цен, при которых $y$ слабо разрешим. Ниже задача ставится на нём.
\end{remark}

\begin{lemma}[Каратеодори]\label{lem:cara}
    Пусть $\mathcal A\subset2^{[T]}$ и $y\in\cone\set{\mathbf 1_S\mid S\in\mathcal A}$. Тогда $y=\sum_{k\in[T]}m_k\mathbf 1_{S_k}$ при некоторых $m_k\ge0$ и $S_k\in\mathcal A$.
\end{lemma}

\begin{proof}
    Среди представлений $y=\sum_{S\in\mathcal A}m_S\mathbf 1_S$ с $m_S\ge0$ возьмём представление с наименьшим числом ненулевых коэффициентов. Если их больше $T$, векторы $\mathbf 1_S$ с $m_S>0$ линейно зависимы в $\R^T$: $\sum_Sg_S\mathbf 1_S=0$, $g\ne0$, и, меняя знак, считаем, что некоторое $g_S>0$. Положим $\theta:=\min_{g_S>0}m_S/g_S$. Коэффициенты $m_S-\theta g_S$ неотрицательны, дают то же $y$ и обращаются в нуль ещё в одном месте. Противоречие с минимальностью.
\end{proof}

\begin{lemma}[гипербола]\label{lem:full}
    Пусть $P:=\set{p_t\in\Rpp[d]}_{t\in[T]}$~--- набор цен, у которого произведение первых двух координат одно и то же: $p^1_tp^2_t=\kappa$ при всех $t\in[T]$. Тогда $\widehat{\Sp}(P)=2^{[T]}$, а если первые координаты $p^1_t$ попарно различны, то и $\Sp(P)=2^{[T]}$. В частности, всякий вектор выпусков $y\in\Rp[T]$ слабо разрешим при ценах $P$.
\end{lemma}

\begin{proof}
    Положим $\xi_t:=\bigl((p^1_t)^{-1},(p^2_t)^{-1},0,\dots,0\bigr)$. Так как $p^2_s/p^2_t=p^1_t/p^1_s$, при $s\ne t$
    \[
        \ip{\xi_t}{p_s-p_t}=\frac{p^1_s}{p^1_t}+\frac{p^1_t}{p^1_s}-2=\frac{(p^1_s-p^1_t)^2}{p^1_sp^1_t}\ge0,
    \]
    причём неравенство строгое при $p^1_s\ne p^1_t$. Поэтому $\xi_t$ служит слабым свидетелем индекса $t$ для всякого $S\ni t$, а при попарно различных первых координатах и свидетелем. Второе утверждение: $y=\sum_ty_t\mathbf 1_{\set t}$.
\end{proof}

\begin{lemma}[коробка]\label{lem:box}
    Пусть $p,q>0$ и $D\ge0$. Тогда $\pi(p,q)=\max\set{q/p,\,p/q}-1$, и $\pi(p,q)\le D$ тогда и только тогда, когда $p/(1+D)\le q\le p(1+D)$. В частности, если $\pi(P,Q)\le D$, то $p^i_t/(1+D)\le q^i_t\le p^i_t(1+D)$ при всех $t\in[T]$, $i\in[d]$.
\end{lemma}

\begin{proof}
    Положим $a:=q/p$. При $a\ge1$ имеем $\abs{a-1}=a-1\ge1-a^{-1}=\abs{a^{-1}-1}$, при $a\le1$~--- симметрично. Поэтому $\pi(p,q)=\max\set{a,a^{-1}}-1$, и $\pi(p,q)\le D$ равносильно $a\le1+D$ и $a^{-1}\le1+D$. Второе утверждение: слагаемые в $\pi(P,Q)$ неотрицательны, поэтому $\pi(p^i_t,q^i_t)\le\pi(P,Q)\le D$.
\end{proof}

\begin{lemma}[оценка сверху]\label{lem:bound}
    Пусть $d\ge2$, $P:=\set{p_t\in\Rpp[d]}_{t\in[T]}$~--- набор цен, $y\in\Rp[T]$, $\kappa>0$ и
    \[
        D(\kappa):=\sum_{t\in[T]}\Bigl(\max\set{\frac{\kappa}{p^1_tp^2_t},\ \frac{p^1_tp^2_t}{\kappa}}-1\Bigr).
    \]
    Тогда найдётся набор цен $Q\in\Rpp[d\times T]$, при котором $y$ слабо разрешим, с $\pi(P,Q)=D(\kappa)$.
\end{lemma}

\begin{proof}
    Положим $q^1_t:=p^1_t$, $q^2_t:=\kappa/p^1_t$ и $q^i_t:=p^i_t$ при $i>2$. Тогда $q^1_tq^2_t=\kappa$, и $y$ слабо разрешим при ценах $Q$ по лемме~\ref{lem:full}. Меняется только вторая координата, и по лемме~\ref{lem:box}
    \[
        \pi(P,Q)=\sum_{t\in[T]}\pi\bigl(p^2_t,\kappa/p^1_t\bigr)=D(\kappa).
    \]
\end{proof}

\begin{lemma}[достижение минимума]\label{lem:min}
    Пусть $d\ge2$, $P:=\set{p_t\in\Rpp[d]}_{t\in[T]}$~--- набор цен и $y\in\Rp[T]$. Тогда
    \[
        \rho(P,y)=\min\set{\pi(P,Q)\mid Q\in\Rpp[d\times T],\ y\ \text{слабо разрешим при ценах}\ Q},
    \]
    и минимум достигается. В частности, $\rho(P,y)\le D(\kappa)$ при всех $\kappa>0$.
\end{lemma}

\begin{proof}
    Обозначим через $\widehat\rho$ точную нижнюю грань в правой части. Множество, по которому она берётся, непусто по лемме~\ref{lem:bound}. Пересечём его с множеством $\set{Q\mid\pi(P,Q)\le c}$ при $c\ge D(\kappa)$. По лемме~\ref{lem:box} пересечение лежит в компактной коробке $\prod_{t,i}[p^i_t/(1+c),\,p^i_t(1+c)]\subset\Rpp[d\times T]$, а замкнуто оно по теореме~\ref{thm:closure} и непрерывности $\pi(P,\cdot)$. Поэтому непрерывная функция $\pi(P,\cdot)$ достигает на нём минимума $\widehat\rho$ в некоторой точке $\bar Q$. Разрешимые цены слабо разрешимы, откуда $\rho(P,y)\ge\widehat\rho$. Обратно, по теореме~\ref{thm:closure} в любой окрестности $\bar Q$ есть цены $Q$, при которых $y$ разрешим, и по непрерывности $\pi(P,Q)$ сколь угодно близко к $\widehat\rho$, откуда $\rho(P,y)\le\widehat\rho$.
\end{proof}

\begin{definition}[Задача о возмущении]\label{def:model}
    Пусть $P:=\set{p_t\in\Rpp[d]}_{t\in[T]}$~--- набор цен и $y\in\Rp[T]$~--- вектор выпусков. Зафиксируем параметр $D>0$, границу возмущения цен, и рассчитаем вспомогательные границы
    \[
        Y:=\max_{t\in[T]}y_t,\qquad M_{st}:=(1+D)\sum_{i\in[d]}(p^i_s+p^i_t).
    \]
    Переменные:
    \begin{itemize}
        \item $a^i_t$~--- множители цен, $q^i_t:=p^i_ta^i_t$;
        \item $b^i_t$ и $g^i_t$~--- обратные множители и возмущения координат;
        \item $z_{k,t}$~--- индикаторы спектров $S_k:=\set{t\mid z_{k,t}=1}$;
        \item $m_k$ и $w_{k,t}$~--- массы и их произведения с индикаторами;
        \item $\xi^i_{k,t}$~--- слабые свидетели.
    \end{itemize}
    \textit{Задачей о возмущении} с параметром $D$ назовём задачу
    \begin{empheq}[left=\empheqlbrace]{align}
        & \sum_{t\in[T],\,i\in[d]}g^i_t\ \to\ \min = v(D), \label{eq:obj}\\
        & g^i_t\ge a^i_t-1,\quad g^i_t\ge b^i_t-1,\quad a^i_tb^i_t\ge1, & &\forall\,t\in[T],\ i\in[d], \label{eq:pert}\\
        & \sum_{k\in[T]}w_{k,t}=y_t, & &\forall\,t\in[T], \label{eq:cone}\\
        & w_{k,t}\le Y z_{k,t},\quad w_{k,t}\le m_k,\quad w_{k,t}\ge m_k-Y(1-z_{k,t}), & &\forall\,k,t\in[T], \label{eq:product}\\
        & \sum_{i\in[d]}\xi^i_{k,t}\,\bigl(p^i_sa^i_s-p^i_ta^i_t\bigr)\ \ge\ -M_{st}\,(1-z_{k,t}+z_{k,s}), & &\forall\,k,s,t\in[T],\ s\ne t, \label{eq:witness}\\
        & a^i_t,\,b^i_t\in[1/(1+D),\,1+D],\quad g^i_t\in[0,D], & &\forall\,t\in[T],\ i\in[d], \label{eq:box}\\
        & \xi^i_{k,t}\ge0,\quad \sum_{i\in[d]}\xi^i_{k,t}=1, & &\forall\,k,t\in[T],\ i\in[d], \label{eq:simplex}\\
        & z_{k,t}\in\set{0,1},\quad m_k\in[0,Y],\quad w_{k,t}\in[0,Y], & &\forall\,k,t\in[T]. \label{eq:box2}
    \end{empheq}
\end{definition}

\begin{theorem}[значение задачи]\label{thm:value}
    Пусть $d\ge2$ и $D\ge D(\kappa)$ при некотором $\kappa>0$. Тогда задача о возмущении с параметром $D$ разрешима, $v(D)=\rho(P,y)$, и для всякого её оптимального решения вектор $y$ слабо разрешим при ценах $Q:=\set{p_t\circ a_t}_{t\in[T]}$, причём $\pi(P,Q)=\rho(P,y)$.
\end{theorem}

\begin{proof}
    Пусть точка $(a,b,g,z,m,w,\xi)$ допустима и $q^i_t:=p^i_ta^i_t$. Из $a^i_tb^i_t\ge1$ и $a^i_t,b^i_t\ge0$ следует $a^i_t>0$ и $b^i_t\ge1/a^i_t$, поэтому $Q\in\Rpp[d\times T]$ и по лемме~\ref{lem:box} $\pi(p^i_t,q^i_t)=\max\set{a^i_t,1/a^i_t}-1\le\max\set{a^i_t,b^i_t}-1\le g^i_t$. Ограничения~\eqref{eq:product} дают $w_{k,t}=m_kz_{k,t}$, поэтому~\eqref{eq:cone} означает $y=\sum_km_k\mathbf 1_{S_k}$. При $t\in S_k$ и $s\notin S_k$ правая часть~\eqref{eq:witness} равна нулю, так что $\ip{\xi_{k,t}}{q_s-q_t}\ge0$ и $\xi_{k,t}$~--- слабый свидетель индекса $t$. Значит, $S_k\in\widehat{\Sp}(Q)$, вектор $y$ слабо разрешим при ценах $Q$, и $\pi(P,Q)\le\sum_{t,i}g^i_t$.
    Обратно, пусть $y$ слабо разрешим при ценах $Q\in\Rpp[d\times T]$ и $\pi(P,Q)\le D$. Положим $a^i_t:=q^i_t/p^i_t$, $b^i_t:=p^i_t/q^i_t$ и $g^i_t:=\pi(p^i_t,q^i_t)$. По лемме~\ref{lem:cara} $y=\sum_{k\in[T]}m_k\mathbf 1_{S_k}$ с $S_k\in\widehat{\Sp}(Q)$, при $S_k=\varnothing$ считаем $m_k=0$. Положим $z_{k,t}:=[t\in S_k]$, $w_{k,t}:=m_kz_{k,t}$, при $t\in S_k$ возьмём в качестве $\xi_{k,t}$ слабого свидетеля индекса $t$ с $\sum_i\xi^i_{k,t}=1$, а при $t\notin S_k$~--- любой вектор с такой суммой, откуда~\eqref{eq:simplex}. Тогда $a^i_tb^i_t=1$ и $g^i_t=\max\set{a^i_t,b^i_t}-1\le\pi(P,Q)\le D$, откуда~\eqref{eq:pert} и~\eqref{eq:box}. Далее, $m_k\le y_t\le Y$ при $t\in S_k$, откуда~\eqref{eq:cone}, \eqref{eq:product} и~\eqref{eq:box2}. В~\eqref{eq:witness} при $t\in S_k$, $s\notin S_k$ правая часть равна нулю, а левая равна $\ip{\xi_{k,t}}{q_s-q_t}\ge0$. В остальных случаях правая часть не больше $-M_{st}$, а левая не меньше $-\sum_i\xi^i_{k,t}(q^i_s+q^i_t)\ge-(1+D)\sum_i(p^i_s+p^i_t)=-M_{st}$. Значение целевой функции равно $\pi(P,Q)$.
    По лемме~\ref{lem:min} минимум $\pi(P,Q)$ по слабо разрешимым $Q$ достигается и равен $\rho(P,y)\le D(\kappa)\le D$. Вторая часть даёт допустимую точку со значением $\rho(P,y)$, первая показывает, что меньшего значения не бывает, и что всякое оптимальное решение даёт слабо разрешимый $Q$ с $\pi(P,Q)\le v(D)=\rho(P,y)$.
\end{proof}

\begin{remark}
    Задача~\eqref{eq:obj}--\eqref{eq:box2}~--- смешанно-целочисленная задача с линейной целевой функцией, бинарными переменными $z$ и билинейными членами $\xi^i_{k,t}a^i_s$ в~\eqref{eq:witness}. Множество $\set{ab\ge1,\ a,b\ge0}$ выпукло (повёрнутый конус второго порядка), так что~\eqref{eq:pert} невыпуклости не добавляет. Нижняя граница в~\eqref{eq:box} следует из $a^i_tb^i_t\ge1$ и $b^i_t\le1+D$, но задана явно: оболочки Маккормика для $\xi^i_{k,t}a^i_s$ строятся по границам переменных. Все переменные ограничены. Такие задачи Gurobi решает до глобального оптимума пространственным методом ветвей и границ с оболочками Маккормика для билинейных членов~\cite{Gurobi}. Величина \texttt{MIPGap} относится к глобальной нижней оценке, поэтому статус \texttt{OPTIMAL} сертифицирует $v(D)=\rho(P,y)$ с точностью до допусков решателя, а при $D\ge D(\kappa)$ статус \texttt{INFEASIBLE} невозможен по теореме~\ref{thm:value}. Единственный параметр $D$ вычисляется по $P$ до запуска, например $D:=D(\kappa)$ при $\kappa$, минимизирующем $D(\kappa)$. Билинейных членов $dT^3$, бинарных переменных $T^2$. Спектры $S_k$ взаимозаменяемы, поэтому полезно добавить $m_1\ge\dots\ge m_T$.
\end{remark}

\begin{remark}
    Подготовка задачи перед запуском решателя.
    \begin{enumerate}
        \item \emph{Нормировка.} Задача инвариантна относительно покоординатного масштабирования $p^i_t\mapsto\lambda_ip^i_t$ с $\lambda_i>0$, общим для всех $t$: свидетели переходят в $\xi^i/\lambda_i$, разрешимость сохраняется при замене $x^i\mapsto x^i/\lambda_i$, а возмущение $\pi$ не меняется. Это $d$ свободных множителей, и ими стоит выровнять координаты, например так, чтобы $\max_tp^i_t=1$ при каждом $i$. Тогда коэффициенты $p^i_sa^i_s$ и константы $M_{st}$ порядка единицы, а разброс внутри каждой координаты не больше, чем он есть в данных.
        \item \emph{Выбор $\kappa$.} Положим $\pi_t:=p^1_tp^2_t$. По $\lambda:=\ln\kappa$ функция $D(\kappa)=\sum_t(e^{\abs{\lambda-\ln\pi_t}}-1)$ выпукла, и её минимум находится за один проход по отсортированным $\pi_{(1)}\le\dots\le\pi_{(T)}$: для разбиения на $k$ нижних и $T-k$ верхних значений стационарная точка равна $\kappa_k:=\bigl(\sum_{j>k}\pi_{(j)}\big/\sum_{j\le k}\pi_{(j)}^{-1}\bigr)^{1/2}$, и минимум достигается в том $\kappa_k$, которое попадает в $[\pi_{(k)},\pi_{(k+1)}]$, либо в точке излома $\pi_{(k)}$, где производная меняет знак. Берём $D:=D(\kappa)$ с этим $\kappa$.
        \item \emph{Начальная точка.} Набор $Q$ из леммы~\ref{lem:bound} вместе с $S_k=\set k$, $m_k=y_k$ и свидетелями $\xi_t$ из леммы~\ref{lem:full}, нормированными на единичную сумму, даёт допустимую точку со значением $D(\kappa)$. Она передаётся решателю как стартовое решение.
        \item \emph{Фиксация и симметрия.} При $y_t=0$ фиксируем $z_{k,t}=0$ для всех $k$. Спектры $S_k$ взаимозаменяемы, поэтому добавляем $m_1\ge\dots\ge m_T$.
        \item \emph{Сжатие $D$.} Коробке достаточно $D\ge\rho(P,y)$. Найдя допустимую точку со значением $c<D$, задачу перезапускаем с $D:=c$ и ею же как стартовой: оптимум не теряется, а границы в~\eqref{eq:box} и константы $M_{st}$ сжимаются, что усиливает оболочки Маккормика.
        \item \emph{Индикаторы.} Ограничение~\eqref{eq:witness} можно записать без $M_{st}$ как индикаторное: $u_{k,s,t}=1\Rightarrow\sum_i\xi^i_{k,t}(p^i_sa^i_s-p^i_ta^i_t)\ge0$ с бинарной $u_{k,s,t}\ge z_{k,t}-z_{k,s}$. Релаксация обычно сильнее, бинарных переменных становится $T^3$, поэтому при больших $T$ сравниваются оба варианта.
    \end{enumerate}
\end{remark}

\begin{remark}
    Оптимальные цены лежат на границе множества разрешимости, поэтому округление найденного решения может нарушить слабую достижимость спектров $S_k$. Точно сертифицируется верхняя оценка: по теореме~\ref{thm:closure} вблизи найденных цен есть рациональные цены $Q'$, при которых $y$ разрешим, их свидетели имеют положительный запас и проверяются в точной арифметике вместе с системой $\sum_km_k\mathbf 1_{S_k}=y$. Тогда $\rho(P,y)\le\pi(P,Q')$, а нижняя оценка $\rho(P,y)$~--- глобальная граница решателя.
\end{remark}

\section{Алгоритмическая постановка}

\subsection{Сведение к конечной задаче}

\begin{theorem}[сведение]\label{thm:reduction}
    Пусть $d\ge2$, $P:=\set{p_t\in\Rpp[d]}_{t\in[T]}$~--- набор цен и $y\in\Rp[T]$. Тогда
    \[
        \rho(P,y)=\min\set{\pi(P,Q)\ \Bigm|\ Q\in\Rpp[d\times T],\ y\in\cone\set{\mathbf 1_S\mid S\in\widehat{\Sp}(Q)}},
    \]
    минимум достигается, а условие под знаком минимума не содержит ни функции $h$, ни меры $\mu$: это существование конечного семейства множеств $S\subset[T]$, масс $m_S\ge0$ с $\sum_Sm_S\mathbf 1_S=y$ и слабых свидетелей $\xi\in\Rp[d]$, подчинённых конечному числу неравенств $\ip{\xi}{q_s-q_t}\ge0$.
\end{theorem}

\begin{proof}
    Это лемма~\ref{lem:min}. Определение~\ref{def:neo-solvability} разрешимости заменено условием $y\in\cone\set{\mathbf 1_S\mid S\in\Sp(Q)}$ по теореме~\ref{thm:solvability}, а точная нижняя грань по разрешимым $Q$~--- минимумом по слабо разрешимым $Q$ по теореме~\ref{thm:closure}. Каратеодори (лемма~\ref{lem:cara}) ограничивает семейство $T$ множествами.
\end{proof}

\begin{remark}
    Исходная задача ставится в функциональном пространстве: по всем $h\in\Phi_d$, всем окрестностям цен и всем абсолютно непрерывным мерам. Теорема~\ref{thm:reduction} оставляет конечный объект: набор цен $Q$ и комбинаторное семейство спектров. Функция $h$ восстанавливается по нему явно (лемма~\ref{lem:patch}), мера~--- по разложению $y=\sum_Sm_S\mathbf 1_S$ (лемма~\ref{lem:signature}). Дальше вся задача решается в этой конечной постановке.
\end{remark}

\subsection{Блокировка и отсечения}

\begin{definition}[Блокировка]\label{def:blocked}
    Пусть $Q:=\set{q_t\in\Rpp[d]}_{t\in[T]}$, $t\in[T]$ и $R\subset[T]$. Индекс $t$ \textit{заблокирован} множеством $R$ при ценах $Q$, если существует ненулевой $\xi\in\Rp[d]$, для которого $\ip{\xi}{q_r-q_t}\ge0$ при всех $r\in R$.
\end{definition}

Если $S\in\widehat{\Sp}(Q)$, $t\in S$ и $R\cap S=\varnothing$, то $t$ заблокирован множеством $R$: подходит слабый свидетель индекса $t$.

\begin{lemma}[одна точка]\label{lem:block1}
    Индекс $t$ заблокирован множеством $\set s$ тогда и только тогда, когда $q^i_t\le q^i_s$ хотя бы при одном $i\in[d]$.
\end{lemma}

\begin{proof}
    Если $q^i_t\le q^i_s$, подходит $\xi:=e_i$. Если $q_s<q_t$ покоординатно, то $\ip{\xi}{q_s-q_t}<0$ при всяком ненулевом $\xi\ge0$.
\end{proof}

\begin{lemma}[отсечение]\label{lem:cut}
    Пусть $y$ слабо разрешим при ценах $Q$, $w\in\R^T$, $\ip{w}{y}<0$, и $C$~--- конечный набор пар $(t,R)$, $t\in[T]$, $R\subset[T]$, такой что всякое $S\subset[T]$ с $\sum_{t\in S}w_t<0$ содержит $t$ и не пересекает $R$ для некоторой пары $(t,R)\in C$. Тогда хотя бы для одной пары $(t,R)\in C$ индекс $t$ заблокирован множеством $R$ при ценах $Q$.
\end{lemma}

\begin{proof}
    Пусть ни одна пара не заблокирована. Возьмём $y=\sum_Sm_S\mathbf 1_S$ с $m_S\ge0$, $S\in\widehat{\Sp}(Q)$. Если $m_S>0$ и $\sum_{t\in S}w_t<0$, то найдётся $(t,R)\in C$ с $t\in S$, $R\cap S=\varnothing$, и слабый свидетель индекса $t$ блокирует $R$. Противоречие. Значит, $\ip{w}{y}=\sum_Sm_S\sum_{t\in S}w_t\ge0$.
\end{proof}

\begin{corollary}[доминирование]\label{cor:dominance}
    Если $y$ слабо разрешим при ценах $Q$ и $q_s<q_t$ покоординатно, то $y_t\le y_s$.
\end{corollary}

\begin{proof}
    Иначе $w:=\mathbf 1_{\set s}-\mathbf 1_{\set t}$ и $C:=\set{(t,\set s)}$ удовлетворяют лемме~\ref{lem:cut}, и $t$ заблокирован $\set s$, что противоречит лемме~\ref{lem:block1}.
\end{proof}

\begin{lemma}[релаксация]\label{lem:relax}
    Пусть условие $\Phi(Q)$ выполнено при всех $Q\in\Rpp[d\times T]$, при которых $y$ слабо разрешим, и $\rho(P,y)\le D$. Тогда $\rho(P,y)\ge\inf\set{\pi(P,Q)\mid\Phi(Q),\ p^i_t/(1+D)\le q^i_t\le p^i_t(1+D)}$.
\end{lemma}

\begin{proof}
    Минимум в теореме~\ref{thm:reduction} достигается в точке $\bar Q$, она удовлетворяет $\Phi$ и по лемме~\ref{lem:box} лежит в коробке.
\end{proof}

\paragraph{Алгоритм отсечений.} Пусть $D\ge D(\kappa)$.
\begin{enumerate}
    \item \emph{Мастер.} Минимизируем $\pi(P,Q)$ по коробке леммы~\ref{lem:box} при накопленных отсечениях «хотя бы одна пара из $C$ заблокирована». Начальные отсечения~--- $C=\set{(t,\set s)}$ для всех пар с $y_t>y_s$ (следствие~\ref{cor:dominance}). По лемме~\ref{lem:relax} значение мастера~--- нижняя оценка $\rho(P,y)$.
    \item \emph{Разделение.} В оптимуме $Q^*$ мастера решаем задачу линейного программирования «$y\in\cone\set{\mathbf 1_S\mid S\in\widehat{\Sp}(Q^*)}$?» генерацией столбцов. Новый столбец~--- множество $S$, на котором максимальна сумма двойственных переменных при условиях $x_t\le\sum_{r\in R}x_r$ для уже найденных строгих покрытий $(t,R)$. Найденное $S$ проверяется линейной программой из замечания после определения~\ref{def:reachable}. Если $y$ в конусе, $Q^*$ оптимален.
    \item \emph{Отсечение.} Иначе двойственная задача даёт $w$ с $\ip{w}{y}<0$ и $\sum_{t\in S}w_t\ge0$ на $\widehat{\Sp}(Q^*)$. Пока есть $S$ с $\sum_{t\in S}w_t<0$, добавляем в $C$ пару $(t,R)$, где $t\in S$ строго покрывает $R\subset[T]\setminus S$ при ценах $Q^*$. Отсечение $C$ добавляется в мастер.
\end{enumerate}

\begin{theorem}[конечность]\label{thm:finite}
    В точной арифметике алгоритм отсечений останавливается за конечное число шагов, и значение мастера на последнем шаге равно $\rho(P,y)$. Если мастер решается с абсолютной точностью $\varepsilon$, то последний оптимум $Q^*$ слабо разрешим и $\pi(P,Q^*)\le\rho(P,y)+\varepsilon$.
\end{theorem}

\begin{proof}
    Все пары нового отсечения строго покрывают при ценах $Q^*$, то есть не заблокированы. Значит, новое отсечение нарушено в $Q^*$, а все прежние в $Q^*$ выполнены. Поэтому отсечения не повторяются. Пар $(t,R)$ конечное число, наборов пар тоже. На последнем шаге $y$ слабо разрешим при $Q^*$, откуда $\rho(P,y)\le\pi(P,Q^*)$, а обратное неравенство (с поправкой $\varepsilon$) даёт лемма~\ref{lem:relax}.
\end{proof}

\begin{remark}
    Отсечение играет роль отсечения Гомори: оно выводится из сертификата линейной программы в текущей точке, отрезает её и верно для всех допустимых точек. Отличие в том, что условие «покрытие разрушено» есть дизъюнкция. Блокировка точкой~--- выбор координаты в лемме~\ref{lem:block1}. Блокировка парой $\set{r,s}$~--- одна из $d$ линейных возможностей $q^i_t\le\min\set{q^i_r,q^i_s}$ или общий свидетель с билинейными членами $\xi^ia^i_r$. Поэтому мастер~--- смешанно-целочисленная задача. Обойти это нельзя: допустимое множество есть объединение невыпуклых частей. Леммы~\ref{lem:block1}, \ref{lem:cut}, \ref{lem:relax} и следствие~\ref{cor:dominance} проверены в Lean (\lean{Algorithms/Cuts.lean}). Из теоремы~\ref{thm:finite} в Lean проверено комбинаторное ядро: отсечения попарно различны и их число не больше числа всех отсечений (\lean{card_le_of_cut_sequence}).
\end{remark}

\subsection{Доминантная релаксация}

\begin{definition}[Доминантная релаксация]\label{def:dom}
    Положим
    \[
        \rho_{\mathrm{dom}}(P,y):=\min\set{\pi(P,Q)\ \Bigm|\ Q\in\Rpp[d\times T],\ \forall\,t,s:\ y_t>y_s\Rightarrow\exists\,i\ q^i_t\le q^i_s}.
    \]
\end{definition}

По следствию~\ref{cor:dominance} $\rho_{\mathrm{dom}}(P,y)\le\rho(P,y)$. Это мастер первого шага алгоритма отсечений.

\begin{theorem}[разложение]\label{thm:split}
    Пусть $E:=\set{(t,s)\mid y_t>y_s}$. Для отображения $\sigma:E\to[d]$ положим
    \[
        \rho_i(\sigma):=\min\set{\sum_{t\in[T]}\pi(p^i_t,q^i_t)\ \Bigm|\ q^i_t\le q^i_s\ \text{при}\ \sigma(t,s)=i}.
    \]
    Тогда $\rho_{\mathrm{dom}}(P,y)=\min_\sigma\sum_{i\in[d]}\rho_i(\sigma)$, и при фиксированном $\sigma$ каждая задача $\rho_i(\sigma)$ решается за полиномиальное время.
\end{theorem}

\begin{proof}
    Условие $\exists\,i\ q^i_t\le q^i_s$ равносильно выбору $\sigma(t,s)$. При фиксированном $\sigma$ ограничения и целевая функция распадаются по координатам. В переменных $u_t:=\ln q^i_t$ задача $\rho_i(\sigma)$~--- минимизация сепарабельной выпуклой функции $\sum_t\bigl(e^{\abs{u_t-\ln p^i_t}}-1\bigr)$ при условиях порядка $u_t\le u_s$, то есть изотонная регрессия с выпуклыми потерями на частичном порядке. Она решается за полиномиальное время~\cite{HochbaumQueyranne}.
\end{proof}

\paragraph{Ветвление по выбору координат.} Теорема~\ref{thm:split} даёт точный алгоритм для $\rho_{\mathrm{dom}}$ без невыпуклых задач. Узел~--- частичное отображение $\sigma$ на подмножестве $E'\subset E$, его оценка~--- сумма $d$ изотонных регрессий с ограничениями только из $E'$. Это релаксация. Если в решении узла ни одна пара $(t,s)\in E$ не даёт $q_s<q_t$, узел допустим. Иначе выбирается такая пара, и узел делится на $d$ потомков по значению $\sigma(t,s)$. Потомок меняет одну координату, поэтому пересчитывается одна изотонная регрессия. Изотонная регрессия решается точно делением по порогу: при каждом $\theta$ множество $\set{t\mid u^*_t>\theta}$~--- наименьшее верхнее множество, минимизирующее $\sum_tf'_t(\theta+)$, то есть задача о замыкании~\cite{HochbaumQueyranne,Picard}.

\begin{remark}
    Если при ценах $Q$ строгие покрытия только одноточечные, то $\widehat{\Sp}(Q)$ состоит из множеств, замкнутых вниз относительно строгого доминирования, и по лемме~\ref{lem:cake} $y\in\cone\set{\mathbf 1_S\mid S\in\widehat{\Sp}(Q)}$ тогда и только тогда, когда $y_t\le y_s$ при $q_s<q_t$. Задача о новом столбце в этом случае~--- задача о замыкании максимального веса и решается одним минимальным разрезом~\cite{Picard}. Строгие покрытия пар дают условия $x_t\le x_r+x_s$, и задача о столбце становится задачей о замыкании в гиперграфе.
\end{remark}

\subsection{Сложность}

\begin{theorem}[NP-трудность]\label{thm:hard}
    Пусть $A_1,\dots,A_d\subset[n]$, $\bigcup_iA_i=[n]$, и $k^*$~--- наименьшее число множеств $A_i$, покрывающих $[n]$. Пусть $0<\delta\le1$ и $M:=(1+\delta)(1+d\delta)$. Рассмотрим индексы $0$ и $(j,c)$, $j\in[n]$, $c\in\set{1,2}$, цены
    \[
        p_0:=\mathbf e,\qquad p^i_{j,c}:=\begin{cases}1+\delta,& j\in A_i,\\ M,& j\notin A_i,\end{cases}
    \]
    и выпуски $y_0:=1$, $y_{j,c}:=2$. Тогда $\rho(P,y)=\rho_{\mathrm{dom}}(P,y)=\delta k^*$. В частности, вычисление $\rho(P,y)$ при $d$, входящем в данные, NP-трудно, и приближение $\rho(P,y)$ с множителем $(1-\varepsilon)\ln n$ так же трудно, как приближение задачи о покрытии~\cite{Feige}.
\end{theorem}

\begin{proof}
    \emph{Оценка сверху.} Пусть $I$~--- покрытие из $k^*$ множеств. Положим $q^i_0:=1+\delta$ при $i\in I$ и $q^i_0:=1$ иначе, остальные цены не меняем. Тогда $\pi(P,Q)=\delta k^*$. Множество $[T]$ слабо достижимо всегда. Множество $U$ всех индексов $(j,c)$ слабо достижимо: для $(j,c)$ возьмём $i\in I$ с $j\in A_i$, тогда $\xi:=e_i$ даёт $\ip{\xi}{q_0-q_{j,c}}=0$. Так как $y=\mathbf 1_{[T]}+\mathbf 1_U$, вектор $y$ слабо разрешим при $Q$.

    \emph{Оценка снизу.} Пусть $Q$ допустима для доминантной релаксации и $\pi(P,Q)\le d\delta$ (иначе нечего доказывать, так как $k^*\le d$). Положим $\sigma_i:=\max\set{q^i_0-1,0}\le d\delta$ и $x_i:=\min\set{\sigma_i/\delta,1}$. Индекс $0$ стоит не меньше $\sum_i\sigma_i\ge\delta\sum_ix_i$. Для каждого $(j,c)$ есть $i$ с $q^i_{j,c}\le q^i_0\le1+\sigma_i$. При $j\notin A_i$ это стоит не меньше $M/(1+d\delta)-1=\delta$. При $j\in A_i$~--- не меньше $(1+\delta)/(1+\sigma_i)-1\ge\delta(1-x_i)/(1+\delta)$. В обоих случаях не меньше $\delta(1-\max_{i:\,j\in A_i}x_i)/(1+\delta)$. Две копии и $\delta\le1$ дают
    \[
        \pi(P,Q)\ge\delta\Bigl(\sum_{i\in[d]}x_i+\sum_{j\in[n]}\bigl(1-\max_{i:\,j\in A_i}x_i\bigr)\Bigr).
    \]
    Пусть $0=\theta_0<\theta_1<\dots<\theta_m$~--- различные значения $x$ вместе с нулём, $I_k:=\set{i\mid x_i\ge\theta_k}$ и $u(I)$~--- число $j$, не покрытых множествами из $I$. По лемме~\ref{lem:cake} скобка равна
    \[
        \sum_{k=1}^m(\theta_k-\theta_{k-1})\bigl(\abs{I_k}+u(I_k)\bigr)+(1-\theta_m)\,n.
    \]
    Каждое непокрытое $j$ покрывается одним множеством, поэтому $\abs{I_k}+u(I_k)\ge k^*$ и $n\ge k^*$. Веса неотрицательны и в сумме равны $1$, откуда $\pi(P,Q)\ge\delta k^*$. Вместе с $\rho_{\mathrm{dom}}\le\rho$ это даёт равенства. Длина записи данных полиномиальна, и $\rho=\delta k^*$ сохраняет отношения приближения.
\end{proof}

\begin{remark}
    Трудность достигается уже на доминантной релаксации и при $d$, растущем вместе с данными. Сложность при фиксированном $d$, в частности при $d=2$, открыта. Равенство $\rho(P,y)=\delta k^*$ проверено в Lean (\lean{minPert_setCover}, \lean{Algorithms/Hardness.lean}).
\end{remark}

\subsection{Верхние оценки}

\begin{lemma}[фиксированное семейство]\label{lem:fixed-family}
    Пусть $y=\sum_{S\in\mathcal F}m_S\mathbf 1_S$ с $m_S>0$, и для каждой пары $t\in S\in\mathcal F$ зафиксирован ненулевой $\xi_{S,t}\in\Rp[d]$. Тогда множество $K$ наборов $Q$ с $\ip{\xi_{S,t}}{q_s-q_t}\ge0$ при всех $t\in S\in\mathcal F$, $s\notin S$~--- выпуклый многогранный конус, содержащий наборы с $q_1=\dots=q_T$. При всех $Q\in K\cap\Rpp[d\times T]$ вектор $y$ слабо разрешим, и $\inf\set{\pi(P,Q)\mid Q\in K\cap\Rpp[d\times T]}$~--- задача выпуклого программирования, значение которой не меньше $\rho(P,y)$.
\end{lemma}

\begin{proof}
    Ограничения линейны и однородны по $Q$, при равных ценах все разности нулевые. Функция $\pi(P,\cdot)$ выпукла, так как $q\mapsto\max\set{q/p,p/q}$ выпукла при $q>0$.
\end{proof}

Чередование «минимум по $Q\in K$ при фиксированных $\xi$» и «выбор $\xi_{S,t}$ с наибольшим запасом при фиксированном $Q$» не увеличивает $\pi$, так как прежние свидетели остаются допустимыми. В алгоритме семейство $\mathcal F$ берётся из оптимума $Q^*$ мастера: спектры, слабо достижимые при $Q^*$, несут как можно большую часть $y$, остаток $y$ дополняется одноточечными множествами. Другая верхняя оценка~--- алгоритм отсечений, в котором блокировка множеством $R$ разрешена только одной координатой: $q^i_t\le q^i_r$ при всех $r\in R$ и некотором $i$. Это сужение мастера, поэтому его значения не являются нижними оценками. Зато каждый его мастер~--- дизъюнкция изотонных ограничений, и он решается тем же ветвлением по выбору координат без невыпуклых задач. Доказательство теоремы~\ref{thm:finite} проходит дословно: алгоритм останавливается в ценах, при которых $y$ слабо разрешим, и их стоимость~--- верхняя оценка $\rho(P,y)$. Оценка не всегда точна. В примере с треугольником она численно равна $1{,}4641$ при $\rho=2/5$, так как оптимум требует свидетеля с двумя ненулевыми координатами.

Найденные цены лежат на границе множества разрешимости. По лемме~\ref{lem:tilt} сдвиг $Q_\delta$ с малым $\delta>0$ даёт цены, при которых $y$ разрешим, со строго положительными запасами свидетелей и $\pi(P,Q_\delta)\to\pi(P,Q)$. Строгие запасы проверяются в точной арифметике после округления.

\subsection{Вычисления}

Алгоритм применён к $51$ набору данных с $T\in[22,26]$ и $d\in\set{2,3,4}$. Нижние оценки~--- значения мастера, верхние проверены точно в рациональной арифметике.
\begin{itemize}
    \item На $45$ наборах зазор между оценками не больше $1{,}1\cdot10^{-6}$ относительно, время от $8$ секунд до $20$ минут. На $42$ из них $\rho=\rho_{\mathrm{dom}}$: достаточно доминантной релаксации и одного–двух шагов.
    \item На $6$ наборах за отведённое время остался зазор от $0{,}4\%$ до $60\%$. Им нужны десятки отсечений по парам с $20$--$36$ парами в каждом, и мастер с билинейными членами решается долго.
    \item Теорема~\ref{thm:split} проверена численно: на всех $51$ наборе сумма $d$ изотонных регрессий при выборе $\sigma$ из оптимума равна $\rho_{\mathrm{dom}}$.
    \item Гипотеза «оптимум доминантной релаксации ломает все доминирования по одной координате» неверна: она выполнена на $24$ наборах из $51$.
    \item Теорема~\ref{thm:hard} проверена численно на $8$ случайных задачах о покрытии ($n\le6$, $d\le5$): значение алгоритма совпало с $\delta k^*$.
    \item Найденные цены дают слабую, а не строгую разрешимость: у части свидетелей нулевой запас. Сдвиг леммы~\ref{lem:tilt} с $\delta=10^{-9}$ даёт строгую разрешимость, проверенную в точной арифметике, ценой прибавки к $\pi$ порядка $10^{-8}$.
    \item Ветвление по выбору координат с точными изотонными регрессиями в узлах находит $\rho_{\mathrm{dom}}$ на всех $51$ наборе за время от $0{,}01$ до $42$ секунд. На $42$ наборах в найденной точке $y$ слабо разрешим, и точная проверка даёт $\rho=\rho_{\mathrm{dom}}$. Здесь $\rho$ получено без невыпуклых задач: нижняя оценка~--- изотонные регрессии и минимальные разрезы, верхняя~--- линейная программа и проверка в рациональной арифметике.
    \item Алгоритм отсечений с блокировкой только одной координатой находит оптимум за секунды там, где точный мастер работает минуты: на двух наборах с $\rho>\rho_{\mathrm{dom}}$ верхняя оценка совпала с нижней за $15$ и $26$ секунд против $4$ и $20$ минут. На открытых наборах он уменьшил зазор с $8{,}8\%$, $42\%$ и $12\%$ до $0{,}4\%$, $1{,}5\%$ и $3{,}0\%$.
    \item Достаточное условие «в оптимуме доминантной релаксации нет строгих покрытий пар» объясняет равенство $\rho=\rho_{\mathrm{dom}}$ лишь на $24$ наборах. На остальных покрытий пар сотни, но они не мешают. Неравенства $y_t\le y_r+y_s$ для строгих покрытий $(t,\set{r,s})$ выполнены почти везде, в том числе там, где $\rho>\rho_{\mathrm{dom}}$. Препятствия к разрешимости коллективные, как во втором примере раздела «Критерий разрешимости».
\end{itemize}

\subsection{Связь с литературой}

Критерий $y\in\cone$ спектров ячеек при \emph{фиксированной} $h$ получен в~\cite{Shananin99}, см. также~\cite{AMS}, где для $h$ вида CES с неизвестным параметром при $d=2$ дан полиномиальный алгоритм. Конечные характеризации существования вогнутых функций по данным восходят к~\cite{Afriat67,Varian84}. Меры согласия данных с моделью и их NP-трудность изучены в~\cite{CDSS14}, там данные исправляются удалением наблюдений или масштабированием бюджетов, а не возмущением цен. Изотонная регрессия с выпуклыми потерями и задача о замыкании~--- \cite{HochbaumQueyranne,Picard}. Отсечения из сертификатов несовместности~--- комбинаторные отсечения Бендерса~\cite{CodatoFischetti,HookerOttosson}.

\subsection{Открытые вопросы}

\begin{enumerate}
    \item Сложность вычисления $\rho$ и $\rho_{\mathrm{dom}}$ при фиксированном $d$, в частности при $d=2$.
    \item Сложность проверки слабой разрешимости при фиксированных ценах, то есть задачи о столбце с покрытиями пар. Ожидается полиномиальность при $d=2$ и трудность при $d\ge3$.
    \item Условия на $(P,y)$, при которых $\rho=\rho_{\mathrm{dom}}$. На реальных данных равенство выполнено на $42$ наборах из $51$.
    \item Усиление отсечений: сейчас каждое отсечение~--- одна дизъюнкция по $20$--$36$ парам, и мастер закрывает его самой дешёвой парой.
    \item Условия, при которых сужение «блокировка одной координатой» точно. На треугольнике оно не точно, на реальных данных совпадает с нижней оценкой везде, где та известна.
    \item Нижняя оценка без невыпуклых задач для блокировки парой. Ослабление «$t$ заблокирован каждым $r\in R$ по отдельности» сохраняет изотонную структуру, но слабо. Ослабления «$q_t$ не доминирует строго $\sum_rw_rq_r$» для конечного набора весов $w$ (лемма \lean{IsBlocked.exists_le_sum}) дают выпуклые задачи в узлах и сходятся медленно.
    \item Полная проверка теоремы~\ref{thm:finite} в Lean: связать ядро \lean{card_le_of_cut_sequence} с алгоритмом.
\end{enumerate}

\begin{thebibliography}{99}
\bibitem{Gurobi} Gurobi Optimization, LLC, Gurobi Optimizer Reference Manual, 2026,
  \url{https://docs.gurobi.com/projects/optimizer/en/current/features/nonlinear.html}.
\bibitem{Shananin99} А.\,А.~Шананин, Непараметрический метод анализа технологической структуры отрасли, Матем. моделирование 11(9) (1999) 116--122.
\bibitem{AMS} A.\,D.~Agaltsov, E.\,G.~Molchanov, A.\,A.~Shananin, Inverse problems in models of resource distribution, J. Geom. Anal. (2018), arXiv:1702.03576.
\bibitem{Afriat67} S.\,N.~Afriat, The construction of utility functions from expenditure data, Internat. Econom. Rev. 8(1) (1967) 67--77.
\bibitem{Varian84} H.\,R.~Varian, The nonparametric approach to production analysis, Econometrica 52(3) (1984) 579--597.
\bibitem{CDSS14} L.~Cherchye, B.~De Rock, B.~Smeulders, F.\,C.\,R.~Spieksma, Goodness-of-fit measures for revealed preference tests: complexity results and algorithms, ACM Trans. Econ. Comput. 2(1) (2014).
\bibitem{HochbaumQueyranne} D.\,S.~Hochbaum, M.~Queyranne, Minimizing a convex cost closure set, SIAM J. Discrete Math. 16(2) (2003) 192--207.
\bibitem{Picard} J.-C.~Picard, Maximal closure of a graph and applications to combinatorial problems, Management Sci. 22(11) (1976) 1268--1272.
\bibitem{Feige} U.~Feige, A threshold of $\ln n$ for approximating set cover, J. ACM 45(4) (1998) 634--652.
\bibitem{CodatoFischetti} G.~Codato, M.~Fischetti, Combinatorial Benders' cuts for mixed-integer linear programming, Oper. Res. 54(4) (2006) 756--766.
\bibitem{HookerOttosson} J.\,N.~Hooker, G.~Ottosson, Logic-based Benders decomposition, Math. Program. 96(1) (2003) 33--60.
\end{thebibliography}

\end{document}
